\section{Results}

\subsection{Three prospectively valid frontier questions}
The bounded series contains three prospectively frozen scientific questions. The inferential unit is one frontier question, not a receipt, controller state, exact-search row, or host call. Two earlier candidate questions were retired before scoring when later-result-oriented work became visible before instrument freeze; they remain in the audit history but are excluded from the valid series rather than silently replaced after outcome access.

\subsection{Case 1: regime characterization}
In the first case, Lane A identifies the responsible layer as representation-regime characterization and selects regime-predicate characterization as primary, with a support-two closure test as a complementary probe. Lane B reaches the corresponding typed responsibility state and selects the same primary computational move while withholding a representation revision whose prerequisite obligation is unresolved.

The contemporaneous relation is agreement on the scored diagnosis and primary move. Later exact regime-characterization and support-two results align with the frozen move coordinates. This is one prospective case, not a probability estimate.

\subsection{Case 2: outside-cone sharpness remains open}
The second case freezes an objective immediately outside a sufficient-cone boundary and asks whether an exact support-two method fails relative to unrestricted dynamic programming. The later analyzer checks 53 frozen rows: three hostile one-object rows, two hostile two-object rows, 24 deterministic random two-object rows and 24 deterministic random three-object rows. It finds zero strict cost gaps in which unrestricted dynamic programming beats the support-two method.

The zero-gap result aligns with both instruments' frozen responsibility and next-move coordinates, but its scientific interpretation is deliberately narrower: it is exact finite-panel evidence and does not enlarge the all-size sufficient cone or prove sharpness.

\subsection{Case 3: agreement without diagnostic correctness}
The third case changes the objective weights and freezes a complete recomputation over 729 ordered one-object instances and 38,760 reorder-quotiented two-object instances, for 39,489 instances in total. The pre-existing finite-domain predicate has zero mismatches with the reweighted incumbent-exact label on this complete declared domain.

Before that outcome existed, both instruments agreed on a responsibility interpretation tied to an objective-scoped boundary and selected a complete reweighted census as the next move. Under the scoring map frozen before the analyzer/result, the later evidence supports a stronger finite-domain structural-invariance interpretation. Consequently the two instruments \emph{agree with each other} on responsibility while both responsibility scores are false; their selected move nevertheless scores aligned for both lanes.

This is the central counterexample of the series: inter-instrument agreement and later diagnostic correctness are distinct variables. The benchmark therefore does not treat consensus as validation.

\subsection{Descriptive series boundary}
All three frozen primary moves are later aligned under their case-specific scoring maps, but the questions are non-random and all originate in one research programme. We report the three rows individually. We do not report ``100\% move accuracy'', a binomial interval, a kappa coefficient, a reliability estimate, calibration, or a generalized success probability.